% Sección 3: Calidad, preparación e integración de factores externos
% evidencia:E04,E05,E10,E26

\section{Calidad, preparación e integración de factores externos}

La confiabilidad de los modelos predictivos y de la reconstrucción del censo depende de resolver anomalías estructurales del registro administrativo antes de entrenar o simular. El equipo ejecutó una auditoría SQL de solo conteos sobre las tablas clínicas principales y formalizó reglas de depuración directamente dentro de las consultas de extracción y del preprocesamiento en Python \parencite{mediflow2026calidad,mediflow2026informe}.

\subsection{Auditoría archivada de calidad y consistencia temporal}

Los resultados archivados de la auditoría del 28 de octubre de 2026 \parencite{mediflow2026calidad} documentan completitud total (0,00\,\% de valores nulos) en las columnas clave de \texttt{person} (\num{912205} filas), \texttt{visit\_occurrence} (\num{1706001} filas) y \texttt{visit\_detail} (\num{3281613} filas). Aunque los identificadores primarios (\texttt{visit\_occurrence\_id} y \texttt{visit\_detail\_id}) no presentan repeticiones, el agrupamiento por clave de evento revela duplicidad administrativa significativa en visitas principales y residual en etapas. Adicionalmente, el cruce con el año de nacimiento detectó edades biológicamente imposibles que distorsionaban el grupo pediátrico cuando no se reclasificaban explícitamente.

Para verificar si la anonimización había desplazado artificialmente las fechas respecto del calendario colombiano real, la auditoría midió el volumen relativo de visitas en festivos entre semana frente a las tres semanas vecinas del mismo día de la semana, desplazando la referencia entre $-21$ y $+21$ días \parencite{mediflow2026calidad}. En Rehabilitación Ambulatoria (citas programadas que cierran en días no laborales), el mínimo global se alcanzó exactamente en el desfase de 0 días con un volumen relativo de 0,0032 (0,32\,\%) en \texttt{visit\_occurrence} y 0,0033 (0,33\,\%) en etapas de \texttt{visit\_detail}, con márgenes de 0,50 y 0,49 frente al segundo menor desplazamiento, respectivamente, además de caídas casi totales en lunes y martes de Carnaval (0,25\,\% y 0,21\,\%) y cierre de fin de año. Aunque Urgencias no discrimina por su operación continua 24/7 (73,95\,\% en festivos), la prueba aporta consistencia calendario-registro suficiente para habilitar variables calendáricas reales, sin constituir una certificación absoluta de exactitud para cada marca horaria individual. Todas estas cifras corresponden a resultados archivados del equipo y no a una nueva consulta SQL ejecutada para la edición del presente informe.

\begin{table}[htbp]
\caption{Reglas de calidad de datos, hallazgos archivados e implementación SQL vigente}
\label{tab:reglas_calidad}
{\singlespacing\small
\begin{tabularx}{\textwidth}{@{}L{0.18\textwidth} L{0.30\textwidth} L{0.24\textwidth} Y@{}}
\toprule
Dimensión & Hallazgo archivado y población & Implementación vigente & Efecto o límite \\
\midrule
Nulos críticos & 0,00\,\% en claves y marcas de \texttt{person}, \texttt{visit\_occurrence} y \texttt{visit\_detail} & No requiere imputación de fechas & Validez de completitud condicionada al extracto \\
Duplicados en visitas & \num{228259} copias sobrantes (13,38\,\%) sobre \num{1706001} filas & \texttt{COUNT(DISTINCT (person\_id, concept, inicio))} en SQL \texttt{04}, \texttt{11}, \texttt{13} & Reduce las llegadas totales a \num{1477742} ($-13{,}4\,\%$) \\
Duplicados en etapas & \num{14081} copias sobrantes (0,43\,\%) sobre \num{3281613} etapas & Conteo agregado por etapa; tasa distinta a visitas & No extrapolar 0,43\,\% como tasa de visitas \\
Edades imposibles & 43 visitas (22 pacientes, 0,0025\,\%): edad $<0$ o $>110$ años & Etiqueta \texttt{invalida} en SQL \texttt{10} y \texttt{13} & Evita inflar el grupo pediátrico \texttt{0-17} \\
Etapas inválidas y huérfanas & 0\,\% sin fin o con fin $<$ inicio; 117 huérfanas globales (91 en \texttt{9201}) & SQL \texttt{14}--\texttt{16}, \texttt{18} excluyen fechas inválidas; \texttt{SQL16} usa \texttt{LEFT JOIN} y conserva huérfanas & Ver contraste poblacional en el apartado de recorte $p_{99{,}9}$ \\
Regularización temporal & Ausencia de filas en horas o días sin movimientos & \texttt{grilla\_regular} rellena con 0 en \texttt{agregados.py} & Evita sesgar al alza los promedios horarios y diarios \\
\bottomrule
\end{tabularx}}
\apanota{Síntesis contrastada entre \textcite{mediflow2026calidad}, \textcite{mediflow2026informe} (sección~3) y los archivos SQL \texttt{04}, \texttt{13}, \texttt{14}, \texttt{15}, \texttt{16} y \texttt{18} de \texttt{src/hun/data/sql/}.}
\end{table}

\subsection{Depuración estadística de estancias extremas (\texorpdfstring{$p_{99{,}9}$}{p99.9}) y regla vigente sobre huérfanas}

La auditoría sin recorte ejecutada en \texttt{23\_duracion\_etapas\_percentiles.sql} reveló colas administrativas extremas que distorsionaban la reconstrucción acumulada del censo horario: los máximos crudos alcanzaron \num{28070,3}\,h en Tratamiento Urgencias, \num{20730,3}\,h en Urgencias y \num{11213,6}\,h en Hospitalización \parencite{mediflow2026saturacion,mediflow2026informe}. Para evitar que registros abiertos o anomalías de cierre inflaran artificialmente la ocupación acumulada, las consultas \texttt{14\_etapas\_entradas\_por\_hora.sql}, \texttt{15\_etapas\_salidas\_por\_hora.sql} y \texttt{16\_duracion\_etapas\_diaria.sql} aplican como filtro superior el percentil $p_{99{,}9}$ calculado por cada tipo de etapa sobre duraciones válidas: \textbf{\num{1176,3}\,h} en Hospitalización (\num{1176,257222}\,h exactas), \textbf{\num{283,7}\,h} en Tratamiento Urgencias, \textbf{\num{17,5}\,h} en Urgencias y \textbf{0,0\,h} en Rehabilitación Ambulatoria \parencite{mediflow2026saturacion}.

En Hospitalización (\texttt{9201}), sobre un universo de \num{108601} etapas válidas, exactamente 107 superan el umbral de \num{1176,3}\,h (todas ellas asociadas a visita padre) y quedan excluidas, conservando \textbf{\num{108494} etapas} incluidas las 91 etapas huérfanas (cuya duración individual es inferior a 720\,h) \parencite{mediflow2026informe}. Al aplicar este filtro sobre el censo horario posterior al calentamiento, los percentiles $p_{50}/p_{75}/p_{90}$ de Hospitalización pasan de 288 / 306 / 319 (máximo 378) a \textbf{283 / 297 / 309} (máximo \textbf{359}), mientras los de Tratamiento Urgencias pasan de 219 / 260 / 300 a \textbf{201 / 239 / 272} \parencite{mediflow2026informe}. Este filtro debe interpretarse estrictamente como un recorte estadístico agregado que puede excluir estancias reales prolongadas; no constituye una adjudicación clínica caso a caso.

El contraste directo del código SQL permite resolver una divergencia documental existente en el repositorio. Mientras algunas frases históricas de \textcite{mediflow2026calidad} recomiendan excluir las etapas huérfanas sin registro en \texttt{visit\_occurrence}, la implementación vigente en \texttt{16\_duracion\_etapas\_diaria.sql} utiliza \texttt{LEFT JOIN visit\_occurrence} bajo el mismo criterio de \texttt{SQL14} y \texttt{SQL15}, conservando la población común de \num{108494} etapas. Si se exigiera visita padre mediante unión interna, la cohorte se reduciría a \num{108403} etapas, definiendo una población distinta \parencite{mediflow2026informe}. En este informe prevalece la regla implementada en el SQL vigente para describir la arquitectura actual del censo alineado, conservando el registro de la población alternativa en el análisis de Ley de Little.

\subsection{Integración externa y protocolo de cuatro filtros}

El módulo independiente \texttt{src/external/} genera la tabla diaria de factores externos (\texttt{data/externos.parquet}) integrando calendario colombiano (librería \texttt{holidays} con festivos, puentes y día posterior al festivo, Carnaval de Barranquilla, Semana Santa, vacaciones escolares y quincenas/primas), meteorología histórica de Open-Meteo para las coordenadas de Soledad y series epidemiológicas semanales de SIVIGILA para dengue e infección respiratoria aguda grave con rezagos anti-fuga de dos y tres semanas (\texttt{lag2s} y \texttt{lag3s}). El archivo de control \texttt{data/externos\_fuentes.json} fechado el 27 de octubre de 2026 certifica estado \texttt{ok} para calendario, clima (del \texttt{2018-01-01} al \texttt{2026-10-21}) y epidemiología; sin embargo, la mera disponibilidad de una fuente no acredita aporte predictivo ni completitud diaria en el horizonte operativo \parencite{mediflow2026informe}.

\begin{table}[htbp]
\caption{Resultados archivados del protocolo de evaluación empírica de factores externos}
\label{tab:factores_externos}
{\singlespacing\small
\begin{tabularx}{\textwidth}{@{}L{0.23\textwidth} l r r r r Y@{}}
\toprule
Grupo evaluado & Decisión & Mediana & Placebo & Pliegues & $p$-signo & Motivo registrado \\
\midrule
\texttt{festivos\_y\_puentes} & Acepta & $+33{,}80\,\%$ & $-1{,}80\,\%$ & 88\,\% & 0,000 & Cumple los 4 criterios \\
\texttt{carnaval\_y\_}\allowbreak\texttt{semana\_santa} & Acepta & $+3{,}49\,\%$ & $+1{,}26\,\%$ & 73\,\% & 0,014 & Cumple los 4 criterios \\
\texttt{vacaciones\_}\allowbreak\texttt{escolares} & Descarta & $+4{,}41\,\%$ & $-0{,}13\,\%$ & 65\,\% & 0,084 & $< 70\,\%$ de pliegues y $p > 0{,}05$ \\
\texttt{quincena\_y\_primas} & Descarta & $+1{,}17\,\%$ & $+0{,}79\,\%$ & 58\,\% & 0,279 & Placebo alto ($> 0{,}5\times$), $< 70\,\%$ pliegues y $p > 0{,}05$ \\
\bottomrule
\end{tabularx}}
\apanota{Tomado de \textcite{mediflow2026externos} (corrida del 18 de noviembre de 2026 sobre 26 pliegues \textit{walk-forward}). Los criterios de aceptación exigen mediana de mejora relativa $\ge 1\,\%$, mejora en $\ge 70\,\%$ de pliegues, prueba unilateral de signo con $p \le 0{,}05$ y mejora del placebo desplazado inferior a la mitad de la mejora real. Los grupos \texttt{clima} y \texttt{epidemiologia} fueron descartados en una corrida inicial anterior \parencite{mediflow2026informe}; la tabla vigente del 18 de noviembre no contiene sus métricas numéricas, por lo que no se reproducen ni interpolan valores para dichos grupos. Las cifras de prueba del placebo de festivos en la nota de \texttt{docs/protocolo\_factores\_externos.md} ($+29{,}6\,\%$, placebo $-0{,}69\,\%$) corresponden a una corrida previa a la actualización del 18 de noviembre.}
\end{table}

Para impedir atribuir a variables exógenas mejoras explicables por estacionalidad o azar, el protocolo somete cada grupo en bloque a cuatro pruebas consecutivas sobre los 26 pliegues de validación temporal: ablación contra un modelo base ya dotado de día de la semana, mes y armónicos anuales; control de estacionalidad; prueba placebo desplazando el factor en el tiempo preservando estructura semanal; y estabilidad mediante porcentaje de pliegues mejorados junto con prueba unilateral de signo \parencite{mediflow2026externos}. La Tabla~\ref{tab:factores_externos} resume los resultados de la corrida del 18 de noviembre de 2026: solo los grupos \texttt{festivos\_y\_puentes} (6 variables) y \texttt{carnaval\_y\_semana\_santa} (5 variables) superan todos los umbrales y son aceptados en el diseño de \texttt{lgbm\_final}.
